\section*{Capítulo \ref{ch:g12:phylogenetic-trees} --- Ler o parentesco: as árvores filogenéticas}

\begin{solution}{exo:g12:phylogenetic-trees:1}
Pontas: as \omterm{def:g10:biodiversity-scales:species}{espécies} comparadas. Nós: os ancestrais comuns mais recentes
hipotéticos delas. Raiz: o ancestral comum de todas elas.
\end{solution}

\begin{solution}{exo:g12:phylogenetic-trees:2}
Os parentes mais próximos do lagarto são o camundongo e o ser humano
juntos (eles compartilham o nó logo abaixo do ramo do lagarto). O
parente mais próximo da rã é o grupo inteiro dos outros três, por igual.
\end{solution}

\begin{solution}{exo:g12:phylogenetic-trees:3}
Um estado que apareceu num ancestral e foi herdado pelos descendentes
dele; compartilhá-lo significa compartilhar esse ancestral. O estado
ancestral era possuído pelo ancestral do grupo inteiro, e assim toda
\omterm{def:g10:biodiversity-scales:species}{espécie} poderia tê-lo conservado: compartilhá-lo não mostra nada sobre
um parentesco mais próximo.
\end{solution}

\begin{solution}{exo:g12:phylogenetic-trees:4}
Um ancestral e todos os descendentes dele. As aves são um \omterm{prop:g12:phylogenetic-trees:rules}{clado}.
``Répteis'' no sentido cotidiano não é: deixa de fora as aves, que
descendem de dentro do grupo.
\end{solution}

\begin{solution}{exo:g12:phylogenetic-trees:5}
O chimpanzé; o nó mais antigo é o que junta o macaco aos primatas
antropoides, 25 milhões de anos atrás.
\end{solution}

\begin{solution}{exo:g12:phylogenetic-trees:6}
Gire em cada nó: a raiz se divide em rã (embaixo) e o resto; depois o
lagarto, depois o camundongo e o ser humano em cima. O camundongo e o
ser humano continuam compartilhando o nó mais recente, o lagarto se
junta a eles em seguida, a rã por último.
\end{solution}

\begin{solution}{exo:g12:phylogenetic-trees:7}
O pombo tem mandíbulas, membros e âmnio: ele entra no \omterm{prop:g12:phylogenetic-trees:rules}{clado} dos
amniotas ao lado do lagarto e do camundongo. As penas, presentes numa
\omterm{def:g10:biodiversity-scales:species}{espécie} só, não definem grupo nenhum aqui; elas definiriam as aves se
mais aves fossem acrescentadas.
\end{solution}

\begin{solution}{exo:g12:phylogenetic-trees:8}
Os pelos, o leite e a mandíbula colocam o golfinho entre os mamíferos;
colocá-lo com o tubarão exigiria que todos os caracteres dos mamíferos
tivessem surgido duas vezes. As nadadeiras e o corpo hidrodinâmico
surgiram separadamente nas duas linhagens: uma convergência.
\end{solution}

\begin{solution}{exo:g12:phylogenetic-trees:9}
$2.2/0.27 \approx 8$ milhões de anos.
\end{solution}

\begin{solution}{exo:g12:phylogenetic-trees:10}
Os dois são pontas do presente; nenhum descende do outro. Correto: os
seres humanos e os chimpanzés são \omterm{def:g10:biodiversity-scales:species}{espécies} irmãs, descendentes de um
ancestral comum, hoje extinto, que viveu cerca de seis milhões de anos
atrás e não era nem um ser humano nem um chimpanzé.
\end{solution}

\begin{solution}{exo:g12:phylogenetic-trees:11}
Sabe-se que a lampreia fica fora do grupo dos \omterm{def:g10:common-ancestry:bodyplan}{vertebrados} com
mandíbula, e assim os estados dela são ancestrais para eles. Com o
camundongo como grupo externo, os pelos, o âmnio, os membros e as
mandíbulas seriam todos lidos como ancestrais e a ausência deles como
derivada: a árvore seria construída de cabeça para baixo.
\end{solution}

\begin{solution}{exo:g12:phylogenetic-trees:12}
A \omterm{def:g12:selection-drift-speciation:population}{população} ancestral carregava vários \omterm{def:g10:universal-dna:gene}{alelos} de cada \omterm{def:g10:universal-dna:gene}{gene}; quando ela
se dividiu duas vezes em rápida sucessão, \omterm{def:g10:universal-dna:gene}{alelos} diferentes foram
separados nos três descendentes em \omterm{def:g10:universal-dna:gene}{genes} diferentes, e assim a história
de um \omterm{def:g10:universal-dna:gene}{gene} agrupa a \omterm{def:g10:biodiversity-scales:species}{espécie} 1 com a 2 e a de outro agrupa a 2 com a 3.
Cada árvore de \omterm{def:g10:universal-dna:gene}{gene} está certa para aquele \omterm{def:g10:universal-dna:gene}{gene}; a árvore das \omterm{def:g10:biodiversity-scales:species}{espécies}
é o veredito majoritário de muitos \omterm{def:g10:universal-dna:gene}{genes}.
\end{solution}

\begin{solution}{exo:g12:phylogenetic-trees:13}
A 1\% por ano, amostras virais colhidas com meses de intervalo já
diferem de modo mensurável: a árvore do vírus resolve quem infectou
quem dentro de uma epidemia, mas ao longo de milhões de anos a
sequência teria sido sobrescrita muitas vezes. A 0.3\% por milhão de
anos, um \omterm{def:g10:universal-dna:gene}{gene} de mamífero muda devagar demais para se ver qualquer
coisa dentro de uma epidemia, mas guarda um registro legível ao longo
de 200 milhões de anos.
\end{solution}

\begin{solution}{exo:g12:phylogenetic-trees:14}
Um fóssil é uma \omterm{def:g10:biodiversity-scales:species}{espécie} própria, com estados derivados próprios, e não
necessariamente a \omterm{def:g12:selection-drift-speciation:population}{população} ancestral; ele fica num ramo lateral perto
do nó em que as penas apareceram. Ele mostra que o nó --- o ancestral
comum das aves e dos parentes reptilianos mais próximos delas --- tem
pelo menos 150 milhões de anos e que as penas precederam a perda dos
dentes e da cauda.
\end{solution}

\begin{solution}{exo:g12:phylogenetic-trees:15}
Uma árvore não tem topo: girar qualquer nó leva os seres humanos para
baixo sem mudar um parentesco. Toda ponta do presente evoluiu pelo
mesmo tempo desde a raiz, e assim nenhuma é mais avançada. E os seres
humanos são uma ponta entre milhões, num ramo dos primatas
antropoides; a forma da árvore registra descendência, não posição.
\end{solution}

\begin{solution}{pb:g12:phylogenetic-trees:1}
\textbf{1.} Esqueleto ósseo (cinco \omterm{def:g10:biodiversity-scales:species}{espécies}): os \omterm{def:g10:common-ancestry:bodyplan}{vertebrados} ósseos; o
tubarão fica de fora.

\textbf{2.} Quatro membros: tetrápodes (salamandra, crocodilo, pombo,
gato). Ovo amniótico: amniotas (crocodilo, pombo, gato). Encaixados:
\omterm{def:g10:common-ancestry:bodyplan}{vertebrados} ósseos $\supset$ tetrápodes $\supset$ amniotas.

\textbf{3.} Não: um estado numa \omterm{def:g10:biodiversity-scales:species}{espécie} só não agrupa nada. Eles seriam
úteis com mais \omterm{def:g10:biodiversity-scales:species}{espécies} compartilhando-os (outras aves, outros
mamíferos, outros répteis).

\textbf{4.} Raiz: tubarão contra o resto (esqueleto ósseo); depois
salmão contra tetrápodes (quatro membros); depois salamandra contra
amniotas (ovo amniótico); entre os amniotas, uma divisão tripla não
resolvida, com as penas no ramo do pombo e os pelos e a abertura no
crânio no do gato.

\textbf{5.} A tabela coloca crocodilo, pombo e gato juntos, mas não
consegue ordená-los: o parente mais próximo do crocodilo fica indeciso.

\textbf{6.} Crocodilo--pombo, 8\%.

\textbf{7.} O gato (15\% de cada um); depois a salamandra (19--20\%);
depois o salmão (24--25\%); depois o tubarão (30--31\%).

\textbf{8.} Elas concordam em todo encaixe; o \omterm{def:g10:universal-dna:information}{DNA} resolve os amniotas:
crocodilo e pombo são irmãos, o gato se junta a eles em seguida.

\textbf{9.} A linhagem do tubarão se separou antes de todas as outras
divergirem entre si; cada uma delas ficou separada do tubarão pelo
mesmo tempo e mostra a mesma distância.

\textbf{10.} O crocodilo é parente mais próximo do pombo que do gato
--- e, pela mesma lógica, que de um lagarto; ``répteis'' sem as aves
não é um \omterm{prop:g12:phylogenetic-trees:rules}{clado}.

\textbf{11.} 19\% ao longo de $2 \times 340$ milhões de anos de ramo:
cerca de 0.028\% por milhão de anos por linhagem, ou seja, 0.056\% por
milhão de anos de separação.

\textbf{12.} Crocodilo--pombo $8/0.056 \approx 140$ milhões de anos; nó
dos amniotas $15/0.056 \approx 270$ milhões de anos.

\textbf{13.} Salmão $24.5/0.056 \approx 440$; tubarão $30.5/0.056
\approx 540$ milhões de anos.

\textbf{14.} Fósseis: 250 e 310; o relógio dá 140 e 270. A ordem está
certa e a data mais antiga fica perto, mas a data crocodilo--aves é
jovem demais: a taxa de um \omterm{def:g10:universal-dna:gene}{gene} só varia entre linhagens, e um \omterm{def:g10:universal-dna:gene}{gene} dá,
na melhor das hipóteses, um relógio grosseiro.

\textbf{15.} A distâncias grandes, muitas posições mudaram mais de uma
vez, e assim as diferenças não se somam mais: o relógio satura. Os nós
antigos são subestimados por esse \omterm{def:g10:universal-dna:gene}{gene}; um \omterm{def:g10:universal-dna:gene}{gene} mais lento, ou muitos
\omterm{def:g10:universal-dna:gene}{genes}, são necessários.

\textbf{16.} (gato, (crocodilo, pombo)) é o menor \omterm{prop:g12:phylogenetic-trees:rules}{clado} que contém o
gato: os amniotas; depois os tetrápodes; depois os \omterm{def:g10:common-ancestry:bodyplan}{vertebrados} ósseos;
depois todas as seis com o tubarão.

\textbf{17.} ``Peixes'' não é um \omterm{prop:g12:phylogenetic-trees:rules}{clado}: o nó que junta tubarão e salmão
é a raiz, cujos descendentes incluem tudo. Amniotas é um \omterm{prop:g12:phylogenetic-trees:rules}{clado}:
crocodilo, pombo, gato e o ancestral comum deles.

\textbf{18.} A salamandra é uma \omterm{def:g10:biodiversity-scales:species}{espécie} do presente, não um
intermediário; ela é o grupo-irmão dos amniotas, compartilhando com
eles o ancestral tetrápode e sem o ovo amniótico, que surgiu depois na
linha deles.

\textbf{19.} Ao lado do ramo entre o nó dos tetrápodes e o nó dos
amniotas: um tetrápode que não é um amniota. A data dele confirma que
os quatro membros já existiam há 360 milhões de anos e que o ovo
amniótico apareceu depois.

\textbf{20.} (tubarão, (salmão, (salamandra, (gato, (crocodilo,
pombo))))); o \omterm{def:g10:universal-dna:information}{DNA} resolveu que crocodilo e pombo são irmãos; a
separação deles data, por esse \omterm{def:g10:universal-dna:gene}{gene}, de cerca de 140 milhões de anos
(os fósseis dizem 250).
\end{solution}
